% ===== §18b: objections and replies =====
\section{Objections and Replies}
\label{p26:sec:objections}

\noindent
Before the close, this section states and answers four objections that the account invites. The aim is not to defend the position against all comers, since the paper is constructive and sets up no adversary, but to make its commitments precise by showing what it does and does not claim. Each objection is given in its strongest form and answered on the account's own terms.

\subsection*{First objection, that the account romanticizes the failure of communication}

The objection runs: by making the limit of language generative, the paper turns a real deprivation, the failure of people to understand one another, into a good, and so counsels acceptance of a distance that causes genuine suffering. It risks telling those who are lonely in their bonds that their loneliness is the condition of love.

The reply distinguishes two things the objection runs together. There is the ordinary failure of communication, in which something sayable is not said, or is said badly, or is misheard, and there is the structural limit, at which something unsayable presses beneath all that is said. The paper makes nothing of the first generative. Ordinary failures of communication are to be repaired, and the account gives no comfort to the withholding of what could be shared or the neglect of what could be understood. What the paper calls generative is only the structural limit, the core that no amount of good communication would reach because it is not of the order of the sayable at all. The distance that is generative is not the distance a better conversation would close; it is the distance that would remain after every good conversation, the room in which the other stays a desiring subject. To confuse the two would indeed be to romanticize a remediable poverty. The paper keeps them apart, and asks only that the irreducible distance, the one no repair reaches, be seen for what it is.

\subsection*{Second objection, that it dissolves the legitimate demand for transparency}

The objection runs: relationships are damaged by concealment, and honesty is a real virtue; an account that prizes the unreachable core seems to license withholding, to dignify the keeping of secrets as fidelity to a mystery.

The reply turns on the difference between the unsayable and the withheld. The core the paper calls unsymbolizable is not a content the subject possesses and declines to share; it is a remainder the subject does not possess either, that presses in them without being available to them as a thing they could disclose or conceal. Withholding concerns the sayable that one chooses not to say, and against withholding the ordinary virtue of honesty stands, undisturbed by this account. The unsymbolizable core is not withheld, because it is not held. To make room for it is not to license secrecy but to acknowledge that even between two people who conceal nothing, who say everything that can be said, a remainder presses in each that neither can disclose because neither possesses it. The demand for transparency, understood as a demand about the sayable, is left entirely in force. What the paper denies is only that transparency, however complete, would reach the core, since the core is not among the things transparency is about.

\subsection*{Third objection, that psychoanalytic premises are doing unearned work}

The objection runs: the argument leans on a particular and contestable theory, the Lacanian account of the subject, the Real, and the drive, and its conclusions are only as secure as those premises, which many do not grant.

The reply concedes the dependence and clarifies its scope. The paper does use the psychoanalytic apparatus, and does not pretend to derive its conclusions from premises everyone shares. But two things soften the objection. First, Part One establishes the six limits phenomenologically, by attention to the dimensions of relational life, before any psychoanalytic term is introduced; a reader who grants the six limits but not the theory has already granted the paper's descriptive core, and the psychoanalytic apparatus of Part Two is offered as the best available account of why the six share one form, not as their only possible ground. Second, the apparatus earns its place by what it explains, the recurrence of one form across six dimensions, the persistence of a core across a life, the death that complete reaching would be, and a reader is entitled to weigh it by that yield rather than by prior allegiance. The account is offered as an explanation to be assessed by its fruit, not as a deduction from axioms held to be self-evident.

\subsection*{Fourth objection, that the paper's own claims symbolize the very gap it calls unsymbolizable}

The objection is the sharpest, and it is reflexive: in stating that the core cannot be symbolized, in giving it the name manque, in setting out claims about it in plain propositions, the paper does to the gap exactly what it says cannot be done, brings it into the signifier. Either the gap can be symbolized, in which case the thesis is false, or it cannot, in which case the paper cannot state its own thesis.

The reply marks a distinction the whole account has relied on, between saying \emph{that} there is a gap and saying \emph{what} the gap contains. The paper does the first and not the second. To state that a remainder is left outside the signifier, to name the place of the remainder, to describe the form of its effects, is to circle the gap, to speak of its edges and its consequences, without delivering its content into the sign. This is exactly the operation the paper attributes to the poem, a circling that keeps the gap open rather than a saying that closes it, and the paper claims for itself no more than this. Its propositions point to the gap and turn around it; they do not fill it. The name manque is not the content of the remainder rendered into a signifier; it is the mark of the place where the content is not. The thesis is therefore consistent with its own performance: a work that circles what it cannot close, and that says so, has not closed it. This is why the paper ends without a synthesis, refusing at the last the closure the objection supposes it must either achieve or fail to achieve. It neither achieves nor fails at closure. It circles, which was the only thing it claimed to do.
